\documentclass[final,5p,times]{elsarticle}

% Compile with XeLaTeX or LuaLaTeX for Korean text.
\usepackage{kotex}
\usepackage{amsmath,amssymb,bm}
\usepackage{graphicx}
\usepackage{booktabs}
\usepackage{multirow}
\usepackage{tabularx}
\usepackage{array}
\usepackage{siunitx}
\usepackage{hyperref}
\usepackage{enumitem}
\usepackage{caption}
\usepackage{subcaption}
\usepackage{placeins}
\usepackage{stfloats}
\usepackage{makecell}

\graphicspath{{./}{./figures/}{./figures/coverage_robustness/}}
\hypersetup{colorlinks=true,linkcolor=blue,citecolor=blue,urlcolor=blue}
\setlist{nosep,leftmargin=1.5em}
\captionsetup{font=small,labelfont=bf}
\journal{Astronomy and Computing}
\sisetup{detect-all}
\emergencystretch=3em
\raggedbottom

% ----------------------------------------------------------------------
% Experiment counts and result macros
% ----------------------------------------------------------------------
\newcommand{\NBackgrounds}{5}
\newcommand{\NBaselineGroups}{3}
\newcommand{\NFluxTiers}{5}
\newcommand{\NControlledRows}{600}
\newcommand{\NCoverageRows}{648}
\newcommand{\NTailRefinedCases}{4}
\newcommand{\NStrictLocalMax}{0}
\newcommand{\NStrictBeamRobust}{0}
\newcommand{\NFullPolyStrict}{0}
\newcommand{\NFrozenOnlyStrict}{0}
\newcommand{\NIntegratedStrict}{0}
\newcommand{\CoverageExploratoryFlags}{16}
\newcommand{\ControlledPTECandidates}{0}
\newcommand{\ControlledMinPTE}{0.0198}
\newcommand{\PTEMinHundred}{1/(100+1)=0.0099}
\newcommand{\PTEmaxStrictRange}{0.022--0.040}
\newcommand{\PTEabsintStrictRange}{0.089--0.969}
\newcommand{\BiasFloorOneEThreeCandidates}{0}
\newcommand{\BiasFloorOneETwoCandidates}{0}
\newcommand{\BiasFloorOneEOneCandidates}{0}
\newcommand{\NominalExpectedFP}{6.48}
\newcommand{\NominalAlpha}{0.01}
\newcommand{\RobustZEpsilon}{10^{-30}}
% Absint floor-hit recheck (N=1000) macros
\newcommand{\NAbsintFloorHitsCoarse}{3}
\newcommand{\NAbsintFloorAboveGate}{2}
\newcommand{\NAbsintRetainedN1000}{2}
\newcommand{\NAbsintBeamRobust}{0}

\begin{document}
\begin{frontmatter}

\title{Starlink 유사 UEMR의 21-cm 지연 스펙트럼 QA에서 국소 최댓값 초과와 윈도우 적분 밴드파워 오염의 불일치 검증}

\author[jnu]{Yujin Kim\fnref{orcid1}}
\ead{europa5218@stu.jejunu.ac.kr}

\author[jnu]{Ji-Won Kang\corref{cor1}}
\ead{jwkang@jejunu.ac.kr}

\affiliation[jnu]{organization={Jeju National University},
                  country={Republic of Korea}}

\cortext[cor1]{Corresponding author}
\fntext[orcid1]{ORCID: \href{https://orcid.org/0009-0006-9745-2960}{0009-0006-9745-2960}}

\begin{abstract}
지구 저궤도(LEO) 위성 군집의 의도치 않은 전자기 방사(UEMR)는
저주파 21-cm 우주론 관측에서 새로운 무선 주파수 간섭(RFI) 요인으로 제기되어 왔다.
그러나 UEMR의 존재 자체가 곧바로 EoR 윈도우 오염을 의미하지는
않으며, 처리된 잔차가 실제 밴드파워 수준에서 통계적ㆍ물리적으로
유의한 편향을 만드는지는 별도의 정량 검증이 필요하다. 본 연구는
Starlink 유사 UEMR 순방향 주입 실험을 통해, HERA 유사 지연 스펙트럼
QA에서 국소 최댓값(local max) 통계량, 경험적 전역 PTE, 절대 윈도우 편향,
윈도우 적분 유의성 및 빔 연산자 강건성이 서로 다른
판정을 줄 수 있음을 보인다.

이를 위해 TLE 기반 순방향 주입과 대응-영가설(matched-null) 보정을 결합하고,
빔 연산자, 가시성 배경, 기선, 플럭스, 스펙트럼 형태학,
위성 다중성을 포함한 통제 그리드를 구성하였다. 또한 native 기선 의존성을 평가하기 위해
2차원 층화추출로 고른 9개 HERA H6C EW-like 기선을 사용하고, 각 기선의 11시간 native
시간축을 10분 연속 청크로 분할하여 정온(quiet)/전형(typical)/스트레스(stress) 대표 빈을
선택한 기선--LST 커버리지 검정을 수행하였다. 통제 그리드에서는 단순한
강건 z(robust-z) 진단량인 $Z_{\rm PS,max}$가 절대 윈도우 편향과 분리되어
국소 플래그를 만들 수 있었고, 전역 PTE 보정 후 강한 오염 후보는
유지되지 않았다.

커버리지 검정에서는 청크 중심 시각이 아니라 10분 청크 전체 시간 표본을 사용하여
위성 노출 메타데이터를 다시 계산하였다. 이 수정 후 거친(coarse) 그리드에서
$PTE_{\rm global,max}<0.01$인 시행은 1건이었으나, 그 시행의 상대 절대 편향은
$10^{-3}$ 보고 기준값보다 작았다. $B_{\rm rel}>10^{-3}$ 조건으로 선별한
\NTailRefinedCases{}개의 경계 근방 사례를 $N_{\rm null}=1000$으로 재평가한 결과,
\NStrictLocalMax{}개의 엄격한 국소 최댓값 후보가 남았다. 윈도우 적분 통계량에서는 고정 빔 한정 바닥값 사례
\NAbsintRetainedN1000{}건이 $N_{\rm null}=1000$에서도 $PTE_{\rm global,absint}<0.01$로 유지되었으나,
짝지은 전체 PolyBeam 시행에서 기각되어 빔-강건 EoR 윈도우 밴드파워 오염 후보는 \NAbsintBeamRobust{}건이다.
핵심 결과는 국소 최댓값 PTE,
절대 편향 기준값, 윈도우 적분 PTE, 전체 빔 강건성을 분리한 다단계 판정
규칙이 필요하다는 점이며, 빈 중심 메타데이터만으로 LST 층을 물리적으로 과해석하면 안 된다는
재현성 점검 결과를 함께 제시한다.
따라서 대응-영가설 앙상블은 UEMR 부재를 탐지하기 위한 영가설이 아니라, 위상 결맞음(phase coherence)에 대한 조건부 스트레스 테스트로 해석되어야 한다.
\end{abstract}

\begin{keyword}
Starlink \sep UEMR \sep HERA \sep EoR \sep delay spectrum \sep RFI \sep matched null \sep empirical PTE \sep QA
\end{keyword}

\end{frontmatter}


\section{서론}
\label{sec:intro}

재이온화 시대(EoR)의 적색편이 21-cm 신호를 측정하려면 극도로 높은 동적 범위가 필요하다.
우주론적 신호는 저주파 가시성(visibility) 데이터에 유입되는 전경, 열잡음, 보정 잔차 및
무선 주파수 간섭(RFI)보다 훨씬 희미하다 \citep{morales2010,deboer2017,abdurashidova2022}.
지연 스펙트럼 및 전경 회피 분석에서 주파수에 대해 매끄러운 전경은 주로 지평선 지연 안쪽의
웻지(wedge)에 집중되는 반면, 지평선 버퍼 바깥쪽의 모드는 EoR 윈도우로 취급된다
\citep{parsons2012delay,liu2014wedge,thyagarajan2015}. 따라서 원자료의 시간--주파수 평면에서는
약해 보이는 구조라도, 처리 연산자를 통과한 뒤 고지연 모드로 재조직되면 운영상 중요한 QA 문제가
될 수 있다.

대규모 LEO 위성 군집은 이 문제를 새롭게 만든다. LOFAR 관측은 110--188 MHz 범위에서
Starlink 위성으로부터 의도치 않은 전자기 방사(UEMR)를 보고했으며 \citep{divruno2023}, 이후
Starlink v2-Mini 관측은 더 낮은 주파수까지 확장되는 빗살(comb) 구조를 포함한 더 강한 저주파 방출을
보고했다 \citep{bassa2024}. 그러나 Starlink 유사 UEMR의 존재와 그것이 21-cm 파워 스펙트럼의
EoR 윈도우 오염으로 전파된다는 명제는 동일하지 않다. 처리된 잔차가 실제 밴드파워
수준에서 통계적ㆍ물리적으로 유의한 편향을 만드는지는 별도로 검증되어야 한다.

본 연구의 중심 질문은 다음과 같다.
\begin{quote}
Starlink 유사 UEMR이 HERA 유사 지연 스펙트럼 QA에서 국소 최댓값 초과를 만들 때, 그것을
실제 윈도우 적분 EoR 윈도우 밴드파워 오염으로 해석할 수 있는가? 또한 국소 최댓값
유의성, 절대 편향, 윈도우 적분 유의성, 전체 빔 강건성은 어떤 순서로
분리되어야 하는가?
\end{quote}

다시 말해, 본 실험은 주입된 UEMR 유사 신호의 결맞음 지연-시간(delay-time) 구조가 선택된 배경, 빔 연산자, 그리고 LST 범위 조건 하에서 위상이 무작위화된 대조군에 비해 QA 통계량을 얼마나 강하게 변화시키는지를 묻는다.

이 질문을 위해 본 논문은 직접적인 역탐지(inverse detection)나 실제 HERA--Starlink 동시 관측을 주장하지 않는다.
대신 TLE 기반 순방향 주입, 대응-영가설 보정, 경험적 PTE, 절대 윈도우 편향 게이트 및
짝지은(paired) 빔 연산자 점검을 결합한 계산적 QA 연구를 수행한다. 본 연구의 결론은 ``Starlink
유사 UEMR이 안전하다'' 또는 ``오염을 검출했다''가 아니라, 국소 최댓값 통계량과 윈도우 적분 밴드파워
오염을 분리하는 다단계 QA 판정 체계가 필요하다는 것이다.

\section{데이터 및 주입 모델}
\label{sec:data}

\subsection{HERA 유사 배경 및 기선 설계}
\label{sec:data:backgrounds}

통제 그리드는 HERA H6C 가시성 크롭을 현실적인 시간--주파수 배경으로 사용하고, 이를 대표
기선 그룹에서 평가한다. 가시성 배경은 HERA $(0,1)$ 기선 산출물로 고정하되, 기하학적
지연 및 지평선 계산에는 통제된 기선 벡터 재정의를 사용한다. 이를 통해 배경 실현(realization)과
지연-지평선 길이를 혼동하지 않는 요인 설계를 구성한다.

HERA H6C 가시성 크롭은 현실적인 시간--주파수 배경으로 사용되지만, 시간 축은 통제된 순방향
주입을 위해 TLE 평가 창으로 재매핑된다. 따라서 본 연구는 동시대 HERA--Starlink 직접 검출이
아니라, 파이프라인 응답과 QA 통계량의 거동을 평가하기 위한 스트레스 실험이다.

\begin{table*}[!t]
\centering
\small
\caption{통제 그리드 및 커버리지 강건성 검정의 요약. 커버리지 검정의 LST 빈 선택은 주입 후 결과값이 아니라 주입 이전 메타데이터만 사용하였다.}
\label{tab:design_overview_revised}
\begin{tabularx}{\textwidth}{>{\raggedright\arraybackslash}p{3.1cm} >{\raggedright\arraybackslash}p{3.2cm} X}
\toprule
실험 축 & 수준 & 분석상 역할 \\
\midrule
통제 그리드 & 배경, 기선, 빔, 플럭스, 형태학, 다중성 & 국소 강건 z, 전역 PTE, 절대 편향의 관계를 점검하는 닫힌 요인 그리드. \\
커버리지 검정 & 9개 native 기선 $\times$ 정온/전형/스트레스 LST 빈 & 길이 계층과 FRF 손실 품질 계층으로 고른 HERA H6C EW-like 기선에서 기선--LST 커버리지 강건성을 점검. \\
꼬리 정밀화 & 경계 근방 4개 사례, $N_{\rm null}=1000$ & 거친 100-영가설 그리드의 PTE 해상도 한계를 보정하고 엄격한 국소 최댓값 후보를 재판정. \\
최종 점검 & $PTE_{\rm global,absint}$, 짝지은 전체 PolyBeam & 국소 최댓값 후보가 윈도우 적분 밴드파워 오염 및 빔-강건 후보로 유지되는지 검정. \\
\bottomrule
\end{tabularx}
\end{table*}

\subsection{순방향 주입}
\label{sec:data:forward}

TLE는 결정론적 감산 모델이 아니라, 물리적으로 가능한 이동 방출원의 시간 종속 방향, 거리, 지연 및
프린지-레이트 궤적을 생성하기 위한 입력값으로 사용된다. 안테나 쌍 $(i,j)$에 대한 근거리 기하학적
지연은 평면파 근사가 아니라 실제 안테나--위성 거리 차이로부터 계산한다.
\begin{equation}
  \tau_{ij}(t)=\frac{|\mathbf{r}_{\rm sat}(t)-\mathbf{r}_{j}|-|\mathbf{r}_{\rm sat}(t)-\mathbf{r}_{i}|}{c}.
\end{equation}
위성 성분의 겉보기 진폭은
\begin{equation}
  A(t,\nu)=S_{\rm ref}\,M_{\rm spec}(\nu)\,B(t,\nu)\,R(t)\,A_{\rm smear}(t,\nu)
\end{equation}
로 정의한다. 여기서 $S_{\rm ref}$는 행 수준 스트레스 스케일, $M_{\rm spec}(\nu)$는 무차원 스펙트럼
템플릿, $B(t,\nu)$는 겉보기 빔 이득, $R(t)$는 거리 감쇠, $A_{\rm smear}(t,\nu)$는 시간 및 주파수 평균에
따른 번짐 계수이다. 주입된 가시성은 다음과 같다.
\begin{equation}
  V_{\rm sat}(t,\nu)=A(t,\nu)\exp[-2\pi i\nu\tau_{\rm sat}(t)].
\end{equation}

빔 연산자는 Gaussian 통제 빔, 150 MHz 고정 PolyBeam, 전체 색수차 PolyBeam으로
분리한다. 본 연구에서 고정 PolyBeam은 빔 근사의 스트레스 사례이며, 전체 PolyBeam은
주파수 의존 빔 연산자에 대한 짝지은 강건성 점검으로 사용한다.

\subsection{플럭스 계층과 스펙트럼 형태학}
\label{sec:data:flux_morphology}

\paragraph{Four-flux integration note.}
$S_{\rm ref}$ is not a calibrated prediction of the HERA-apparent flux of a
particular Starlink transit, but an order-of-magnitude stress-scaling
parameter. Published LOFAR measurements place first-generation broadband UEMR
at roughly 0.1--10 Jy and second-generation v2-Mini broadband UEMR at roughly
2--100 Jy in two 8-MHz windows near 120 and 161 MHz. We therefore use 10 and
30 Jy as literature-scale low/high anchors. The 300 and 1000 Jy levels are
deliberate stress tiers used to probe the dynamic range and failure boundary
of the QA operator; they are not occurrence-rate or calibrated-flux claims.
Beam coupling, slant range, polarization, off-axis response, and emission
directionality are absorbed into the interpretation of $S_{\rm ref}$ as a
stress coordinate rather than a direct observable.

This four-flux design defines the next full-catalog coverage run:
$S_{\rm ref}\in\{10,30,300,1000\}$ Jy. It expands the paired coverage grid
from 648 rows to 1296 rows. Candidate counts, tail refinement, paired-beam
audits, multiplicity bootstrap summaries, ANOVA tables, and candidate-level
calibration audits must therefore be regenerated from the four-flux run. The
existing 648-row result is retained only as the historical two-flux analysis
and is not relabelled as a four-flux result.

$S_{\rm ref}$는 특정 Starlink 통과에 대해 보정된 HERA의 실제 겉보기 플럭스를 예측하는 값이 아니라,
문헌에 보고된 저주파 Starlink UEMR의 크기 범위를 기준으로 파이프라인 응답을 시험하기 위한 주입
척도이다. 본 연구는 낮은 계층과 스트레스 계층을 모두 포함하되, 절대 검출 임계값을 제시하지 않는다.

스펙트럼 형태학은 매끄러운 광대역, 버스트성 광대역, 협대역 선형, kHz-comb 뱅크를 중심으로
구성한다. 초기 분석에서 비매끄러운 형태학은 단순 영가설 또는 특정 연산자 설정에서
국소 QA 플래그를 만들 수 있었으나, 본 논문의 최종 해석은 형태학 자체의 보편적 위험성이 아니라
대응-영가설 보정, 빔 연산자, PTE 종류 및 윈도우 적분 점검에 따라 판정이 달라진다는
점에 둔다.

\section{QA 지표 및 다단계 판정 기준}
\label{sec:metrics}

본 연구의 분석 대상은 전체 HERA 생산 파이프라인의 최종 우주론 산출물이 아니라, 고정된 HERA 유사
가시성 배경 위에 주입된 Starlink 유사 UEMR이 지연-윈도우 QA 지표에 남기는 상대적 변화이다.
주입 전후의 처리 차이는 다음과 같이 둔다.
\begin{equation}
  V_{\rm int}
  =
  P_{\rm QA}(V_{\rm bg}+V_{\rm sat})
  -
  P_{\rm QA}(V_{\rm bg}).
\end{equation}
여기서 $P_{\rm QA}$는 지연 변환, EoR 윈도우 마스크 적용 및 윈도우 통계량 계산까지의 공통 QA 절차를
의미한다.

버퍼가 적용된 EoR 윈도우 프록시는
\begin{equation}
  \tau_{\rm win}=\tau_{\rm hor}+100~{\rm ns}, \qquad
  W_{\rm win}=\{\tau:\ |\tau|\ge \tau_{\rm win}\}
\end{equation}
로 정의한다. 절대 윈도우 편향은
\begin{equation}
  B_{\rm rel}
  =
  \frac{
  \sum_{\tau\in W_{\rm win}}
  \left|
  |\widetilde{V}_{\rm inj}(\tau)|^2-
  |\widetilde{V}_{\rm bg}(\tau)|^2
  \right|
  }{
  \sum_{\tau\in W_{\rm win}}
  |\widetilde{V}_{\rm bg}(\tau)|^2
  }.
\end{equation}
이는 보정된 우주론적 밴드파워 추정량이 아니라 물리적 진폭 게이트로 사용되는 상대 편향 지표이다.

\subsection{대응-영가설 및 경험적 PTE}
\label{sec:metrics:pte}

각 시행에 대해 대응-영가설 앙상블은 주입된 시간--주파수 지지 영역과 진폭 포락선을 보존하되,
위상을 무작위화하여 구성한다.
\begin{equation}
  V_{\rm sat,null}^{(r)}(t,\nu)=|V_{\rm sat}(t,\nu)|\exp[i\phi_r].
\end{equation}
여기서 $\phi_r$의 구체적 의존성은 실험에 따라 다르다. 통제 그리드(matched-null 밴드파워 전파 실험)에서는
$\phi_r(t)$, 즉 적분 시간당 독립적으로 추출된 하나의 위상을 사용하므로 주파수 방향 스펙트럼 결맞음이
보존되지만 시간 방향 fringe-rate 결맞음이 파괴된다. 커버리지 강건성 검정에서는 위성별 전역 스칼라 위상
$\phi_r^{(s)}$를 사용하여 각 위성의 지연--시간 코히런트 구조를 보존한 채 위성 간 상대 위상만 무작위화한다.
단일 위성 커버리지 시행에서는 배경--위성 교차항 위상만 달라지며 위성 자체의 지연 스파이크 구조는 영가설에도
그대로 남는다. 다중 위성 커버리지 시행에서는 위성 간 코히런스 증폭(coherence amplification)이 영가설에서
무작위화되어, 관측값과 영가설의 차이는 위성 간 상대 위상 결맞음에 의한 구조 차이를 반영한다.

이 영가설의 의미는 엄밀히 제한된다. 본 대응-영가설은 ``UEMR 신호가 없다''는 영가설이 아니라,
동일한 지지 영역과 진폭 포락선을 갖는 위상 무작위화 대조군이다. 통제 그리드에서는 시간 방향
위상 스크램블로 fringe-rate 조직화를 제거한 대조군이며, 커버리지 검정에서는 위성별 전역 위상
회전으로 위성 간 코히런스를 제거한 대조군이다. 두 경우 모두, 개별 위성의 지연 영역 에너지 배분은
영가설에서도 변하지 않는다. 따라서 실제 주입 신호에서 국소 통계량이 상승하는 것은 위성 주입의
위상 결맞음 또는 위성 간 코히런스 증폭이 무작위화 대조군보다 강하다는 사실을 반영할 수 있다.
본 연구는 이 차이를 오염 검출로 곧바로 해석하지 않고, 결맞은 구조가 QA 통계량을 얼마나 민감하게
움직이는지를 보는 조건부 스트레스 테스트로 해석한다.

후속(downstream) 프록시 섭동은
\begin{equation}
  \Delta\widehat{P}_{\rm UEMR}(k)
  =
  \widehat{P}_{\rm inj}(k)-\widehat{P}_{\rm bg}(k)
\end{equation}
로 정의하고, 영가설 실현에도 동일한 추정량을 적용한다.

국소 강건 z 진단량은
\begin{equation}
  Z_{\rm PS}(k)=
  \frac{
  \Delta\widehat{P}_{\rm UEMR}(k)
  -
  \mathrm{median}_{r}[\Delta\widehat{P}^{(r)}_{\rm null}(k)]
  }{
  \hat{\sigma}_{r}(k)
  },
\end{equation}
\begin{equation}
  \hat{\sigma}_{r}(k)=
  \begin{cases}
    1.4826\,\mathrm{MAD}_{r}[\Delta\widehat{P}^{(r)}_{\rm null}(k)] & \text{if } > \RobustZEpsilon, \\
    \mathrm{std}_{r}[\Delta\widehat{P}^{(r)}_{\rm null}(k)]          & \text{if MAD } \le \RobustZEpsilon \text{ and std} > \RobustZEpsilon, \\
    \mathrm{NaN}                                                       & \text{otherwise}.
  \end{cases}
\end{equation}
본 연구에서 $Z_{\rm PS,max}$는 검출량이 아니라 국소 탐색적 진단량이다.
\begin{equation}
  Z_{\rm PS,max}=\max_{k\in W_{\rm win}} Z_{\rm PS}(k).
\end{equation}
이 대체 규칙은 영가설 MAD가 수치적으로 0에 수렴하는 퇴화 사례를 처리하기 위한 안전장치이다.
중앙값/MAD 기반 강건 정규화는 극단값에 둔감한 척도 추정량으로 널리 사용된다 \citep{hampel1986,huber1981,rousseeuw1993}.
영가설 MAD가 작아지는 조건에서는 작은 절대 섭동도 큰 강건 z로 증폭될 수 있으므로,
$Z_{\rm PS,max}$ 단독으로 오염을 판정하지 않는다.

경험적 PTE는 다음과 같이 계산한다.
\begin{equation}
  PTE(T_{\rm obs})
  =
  \frac{1+\#\{T_{\rm null}^{(r)}\ge T_{\rm obs}\}}{N_{\rm null}+1}.
\end{equation}
본 연구는 두 종류의 PTE를 구분한다.
\begin{enumerate}
  \item $PTE_{\rm global,max}$: EoR 윈도우 내부 최댓값 통계량에 대한 전역 PTE. 국소 최댓값 초과의 통계적 희귀성을 평가한다.
  \item $PTE_{\rm global,absint}$: 윈도우 적분 절대 통계량에 대한 전역 PTE. 윈도우 적분 밴드파워 수준의 초과 여부를 평가한다.
\end{enumerate}
$PTE_{\rm global,max}<0.01$은 엄격한 국소 최댓값 후보 기준으로 사용하지만, 이 기준만으로
윈도우 적분 오염을 주장하지 않는다. 경험적 PTE는 유한한 영가설 앙상블에서 가능한 해상도가
$1/(N_{\rm null}+1)$로 제한되므로, 엄격한 임계값 근처에서는 꼬리 해상도 점검이 필요하다
\citep{davison1997,north2002}. 또한 다수의 시행에 동일한 임계값을 적용하면 다중 탐색(look-elsewhere) 효과가 발생하므로,
시행 수준 PTE는 그리드 전체의 검출 확률로 직접 해석하지 않는다 \citep{westfall1993}.
본 연구는 $PTE_{\rm global,max}$와 $PTE_{\rm global,absint}$에 동일한 0.01 보고 임계값을 사용하지만,
두 통계량의 의미를 동일시하지 않는다. 전자는 국소 최댓값 초과의 꼬리 사건을, 후자는 윈도우 적분
초과를 평가한다. 같은 숫자 임계값은 표기상 일관된 선별 수준일 뿐이며, 윈도우 적분 오염 주장은
반드시 윈도우 적분 PTE와 빔 강건성 점검을 별도로 통과할 때에만 허용한다.

\subsection{최종 QA 판정 계층}
\label{sec:metrics:hierarchy}

본 연구는 다음의 다단계 판정 구조를 사용한다.
\begin{enumerate}
  \item $Z_{\rm PS,max}$ 또는 $PTE_{\rm global,max}<0.05$: 탐색적/국소 플래그.
  \item $PTE_{\rm global,max}<0.01$: 엄격한 국소 최댓값 통계량 후보.
  \item $B_{\rm rel}>B_{\rm floor}$: 절대 편향 진폭 게이트.
  \item $PTE_{\rm global,absint}<0.01$: 윈도우 적분 후보.
  \item 짝지은 전체 PolyBeam 시행에서도 엄격 기준 유지: 빔-강건 후보.
\end{enumerate}
따라서 코드상 편향 게이트를 통과한 시행이라도, $PTE_{\rm global,absint}$와 전체 PolyBeam 강건성을
통과하지 않으면 윈도우 적분 밴드파워 오염으로 해석하지 않는다. 본 연구는 민감도 보고를 위해
$B_{\rm floor}\in\{10^{-3},10^{-2},10^{-1}\}$를 함께 제시하고, 본 실험의 요약 표에서는 $B_{\rm floor}=10^{-2}$를
운영적 보고 수준으로 사용한다. 이 값은 전체 윈도우 적분 섭동이 배경 윈도우 파워의 1\%를
넘는 경우만 다음 단계 검토 대상으로 남기는 명시적 진폭 게이트이다. 그러나 이는 열잡음 공분산,
보정 잔차 예산, 또는 특정 HERA 상한 분석의 계통오차 하한에 고정한 보편 권고값이 아니다.
실제 생산 수준 QA에서는 해당 관측의 잡음/계통오차 예산에 의해 $B_{\rm floor}$를 사전에 정해야 한다.

\section{실험 설계}
\label{sec:experiments}

\subsection{통제 그리드}
\label{sec:experiments:controlled}

통제 그리드는 형태학, 빔, 플럭스, 다중성, 기선, 배경 조건을 교차시켜
$Z_{\rm PS,max}$, 경험적 PTE, 절대 편향의 관계를 평가한다. 이 단계의 목적은 오염을
확정하는 것이 아니라, 국소 강건 z 플래그와 전역 PTE 보정이 어떻게 달라지는지 확인하는 것이다.

\subsection{Native 9-baseline robustness 및 11시간 기선--LST 커버리지 검정}
\label{sec:experiments:coverage}

커버리지 검정의 첫 번째 목적은 단일 0\_1 기선 크롭에 의존하지 않고, native HERA H6C EW-like
기선에서 결과가 어떻게 변하는지 평가하는 것이다 \citep{deboer2017,abdurashidova2022}. 이를 위해
9개 native 기선을 2차원 층화추출로
선택하였다. 첫 번째 축은 기선 길이 계층이며, 이는 delay horizon
$\tau_{\rm hor}=|\mathbf{b}|/c$를 결정한다 \citep{parsons2012delay,liu2014wedge}. 두 번째 축은
fringe-rate filter(FRF) 손실 품질 계층이다. FRF는 fringe-rate 공간에서 visibility를 시간
필터링하는 연산자이며 \citep{parsons2015frf,garsden2024frf}, 이동 송신원 신호의 잔류 power는 기선
기하와 위성 궤도에 따라 달라진다. 각 기선의 FRF 손실 계수, 즉 위성 신호 power 중 FRF 적용 후
남는 비율을 대표 LEO 궤도에서 계산한 뒤, 같은 길이 계층 안에서 low/mid/high 등급으로 나누었다.
FRF 손실이 낮은 기선은 이동 송신원 신호에 더 민감하므로, 각 길이 계층 안에서 보수적, 중간,
낙관적 시나리오를 함께 대표하도록 하였다. 이 설계로 길이 계층과 FRF 손실 품질 계층의 조합에서
총 9개 native 기선을 사용하였다.

두 번째 목적은 각 native 기선 안에서 위성 노출과 배경 상태가 다른 시간 구간을 비교하는 것이다.
이를 위해 각 UVH5 기선 파일의 native 시간축을 10분에 해당하는
연속 표본 묶음으로 잘랐다. 코드에서는 native cadence로부터 $n_{\rm per\,bin}$을 계산하고,
$i$번째 빈을 시간 인덱스 $[i\,n_{\rm per\,bin},(i+1)n_{\rm per\,bin})$로 정의한다. 이후 위성 기하
계산을 위해 각 빈의 내부 cadence를 보존한 채 시작 시각만 2026년 1월 1일 UTC의 대응 10분
구간으로 재매핑하였다. 따라서 본 논문에서 ``LST 빈''이라고 부르는 항목은 새로 등간격 LST 좌표를
보간해 만든 빈이 아니라, native 관측 시간 표본의 10분 연속 청크에 붙인 분석 라벨이다.

각 crop은 zenith 부근의 peak response를 포착하도록 설계되었지만, 위성 pass의 모든 입사각과
fringe-rate evolution을 연속적으로 추적하는 장시간 주입은 아니다. 본 검정은 zenith 부근의 TLE
ensemble, native 기선 층화, 그리고 10분 청크별 위성 노출 재계산으로 이 제한을 부분적으로 보완하는
robustness assessment이다.

9개 기선 각각에 대해 주입 이전 메타데이터만을 사용하여 정온, 전형, 스트레스 빈을 하나씩 선택하였다.
구체적으로 위성 노출 프록시 $E_{\rm sat}$, 플래그 비율 $f_{\rm flag}$, 그리고 가능한 경우 영가설 전용
윈도우 MAD 프록시를 계산하였다. 정온 빈은 낮은 위성 노출 후보 중 플래그 비율이 해당 기선의
중앙값에 가까운 빈으로 정의하였다. 여기서 $E_{\rm sat}$는 빈 중심 시각 하나가 아니라 해당 10분 청크의
모든 시간 표본에서 계산한 빔 가중 위성 노출의 빈 평균 프록시이다. 중심 시각 노출은 이력(provenance)
열로 보존하지만 층화에는 사용하지 않는다. 전형 빈은 강건 척도화된 $(E_{\rm sat}, f_{\rm flag}, {\rm nullMAD})$
공간에서 중앙값 벡터에 가장 가까운 빈으로 선택하였다. 스트레스 빈은
\begin{equation}
  R_{\rm pre}=\frac{E_{\rm sat}}{{\rm nullMAD}_{\rm win}+\epsilon}
\end{equation}
가 가장 큰 빈으로 선택하되, 같은 빈이 이미 선택된 경우 다음 순위 빈을 사용하였다. 중요한 점은 이 층화가
$Z_{\rm PS,max}$, PTE, $B_{\rm rel}$ 등 주입 후 결과값을 사용하지 않았다는 것이다.
이로써 27개 기선--LST 커버리지 셀을 구성하고, 형태학, 플럭스, 다중성, 빔 연산자를
교차하여 \NCoverageRows{}개 축소 그리드 시행을 실행하였다. 이 셀들은 독립적인 우주론 표본이 아니라,
기선 기하와 관측 상태에 대한 커버리지 탐침으로 해석한다.

커버리지 검정에서 ``가시 위성''은 배경 청크의 시간 표본 전체에서 최고도(peak altitude)가
$70^\circ$ 이상인 위성으로 정의한다. TLE 파일에서는 순서상 처음 1200개 고유 기록만
로드하였으므로, 이 TLE 부분집합이 전체 Jan 2026 Starlink 군집 전체를 대표하지 않을 수 있다.
다중 위성 주입에서는 고도 순위 상위 최대 12개 위성을 선택하였다.

\subsection{Synthetic zenith-angle sensitivity analysis}
\label{sec:experiments:za_sweep}

To isolate geometric sensitivity from the heterogeneous TLE ensemble, we
perform an auxiliary synthetic-orbit sweep at azimuth 90 degrees and altitude
550 km. The peak zenith angle is varied over
$5^\circ,15^\circ,30^\circ,45^\circ,60^\circ,70^\circ,$ and $78^\circ$.
Three representative east-west baselines (14.6, 140.4, and 207.3 m), five
background realizations, and two beam conditions are crossed. The primary beam
condition is the full-chromatic HERA H2C CST PolyBeam; $B=1$ is retained only
as a geometry-only diagnostic comparator and is not used as a physical
candidate gate. Holding $S_{\rm ref}=300$ Jy fixed isolates the angular
response, yielding $5\times3\times7\times2=210$ cases. Every case is evaluated
with the same $PTE_{\rm global,max}$, $PTE_{\rm global,absint}$, and
$B_{\rm rel}$ hierarchy used in the main analysis.

The older 210-case zenith-angle sweep used a null-p95 excess metric rather
than the present PTE/$B_{\rm rel}$ candidate hierarchy. Its design can be
reused, but its numerical exceedance counts are not carried into the current
paper unless the cases are rerun under the current operator and null
definition.

\subsection{꼬리 해상도 정밀화}
\label{sec:experiments:tail}

초기 커버리지 그리드의 영가설 앙상블은 $N_{\rm null}=100$이었으므로 가능한 최소 경험적 PTE는
$\PTEMinHundred$이었다. 엄격한 임계값 주변에서 최종 후보 수를 안정적으로 결정하기 위해,
$PTE_{\rm global,max}\le0.03$ 또는 $Z_{\rm PS,max}>3$ 근방이면서 $B_{\rm rel}>10^{-3}$인 경계 근방 \NTailRefinedCases{}개
사례를 $N_{\rm null}=1000$으로 재계산하였다. 이 정밀화 결과를 최종 국소 최댓값 후보 판정에 사용한다.

추가로, 거친 그리드에서 $PTE_{\rm global,max}$ 기준이 아닌 $PTE_{\rm global,absint}$가 $\PTEMinHundred$
바닥값에 도달한 사례를 별도로 확인하였다. 이 사례들은 $PTE_{\rm global,max}\approx0.7$이어서 위의 꼬리
정밀화 선별 조건에서 제외되었다. 해당 사례는 모두 동일 셀
($11\_10$, LST 빈 16, 정온, 고정 PolyBeam, 다중)에 속하며 \NAbsintFloorHitsCoarse{}건이다. 이 중
$B_{\rm rel}>10^{-3}$ 조건을 만족하는 \NAbsintFloorAboveGate{}건을 $N_{\rm null}=1000$으로 재검증하였다.
재검증 결과는 \ref{sec:results:absint}절에서 서술한다.

\section{결과}
\label{sec:results}

\subsection{통제 그리드: 국소 강건 z는 오염 판정 기준이 아니다}
\label{sec:results:controlled}

통제 그리드에서 $Z_{\rm PS,max}$ 기반 국소 플래그는 일부 발생했지만, 절대 윈도우 편향과의
상관은 약했고 전역 PTE 기준의 강한 오염 후보로 유지되지 않았다. 특히
$PTE_{\rm global,max}<0.01$을 만족한 시행은 \ControlledPTECandidates{}건이었고, 최소 PTE도
약 \ControlledMinPTE{} 수준에 머물렀다. 이는 $Z_{\rm PS,max}>3$와 같은 국소 강건 z 임계값이
윈도우 내 최댓값 통계량과 영가설 MAD 민감도의 영향을 받기 때문에 오염 검출 기준으로 직접
승격될 수 없음을 보여준다.

\begin{figure*}[!t]
\centering
\includegraphics[width=0.82\textwidth]{R2_z_vs_pte_global.png}
\caption{$Z_{\rm PS,max}$와 경험적 $PTE_{\rm global,max}$의 관계. 수직선은 $Z_{\rm PS,max}=3$, 수평선은 $PTE=0.05$ 및 $PTE=0.01$을 표시한다. 국소 강건 z 플래그는 전역 PTE 기준의 엄격한 오염 후보와 동일하지 않다.}
\label{fig:z-vs-pte}
\end{figure*}

\subsection{커버리지 검정: 빈 전체 메타데이터 수정 후 엄격한 후보는 남지 않는다}
\label{sec:results:coverage}

11시간 기선--LST 커버리지 검정에서는 거친 100-영가설 그리드에서 \CoverageExploratoryFlags{}개의
탐색적 플래그가 나타났다. 초기 점검에서 $(0\_1,\ \mathrm{LST\ bin}\ 36)$ 셀은 10분 청크의 중심
시각 위성 노출 프록시에서는 가시 위성이 0개였지만, 같은 청크의 모든 native 시간 표본을 따라
위성 고도와 빔 응답을 평가하면 4개의 가시 위성과 고고도 통과를 포함하였다. 따라서 본 공개판에서는
LST 빈 선택 메타데이터를 빈 중심 프록시가 아니라 빈 전체 시간 표본 기준으로 재계산하였다. 이 사례는
오류를 드러낸 진단 예시이며, 수정된 메타데이터 산출과 빈 재선택은 0\_1 하나가 아니라 9개 기선
전체에 동일하게 적용하였다. 이 수정 후
$(0\_1,\mathrm{bin}\ 36)$ 셀은 선택된 전형 셀에서 제외되었고, $0\_1$ 기선의 선택된 전형
셀은 LST 빈 44로 바뀌었다.

9개 기선 각각의 정온/전형/스트레스 셀, 즉 27개 기선--빈 셀의 새 선택으로
\NCoverageRows{}개 커버리지 시행을 전부 재실행한 결과, 거친 $N_{\rm null}=100$ 그리드에서
$PTE_{\rm global,max}<0.01$인 시행은 1건이었다. 그러나 이 시행은
$B_{\rm rel}=5.53\times10^{-4}$로 $10^{-3}$ 보고 기준값을 넘지 못한다. $B_{\rm rel}>10^{-3}$ 및
경계 근방 조건으로 \NTailRefinedCases{}건을 선별하여 $N_{\rm null}=1000$으로 꼬리 정밀화를 수행하면,
\NStrictLocalMax{}건이 $PTE_{\rm global,max}<0.01$ 기준에 남는다. 따라서 수정된 커버리지 결과의 결론은
``단일 셀에 국한된 엄격 반응''이 아니라 ``빈 전체 메타데이터 수정 후 꼬리 정밀화된 엄격한
국소 최댓값 후보 없음''이다.

\begin{table*}[!t]
\centering
\small
\caption{수정된 선택 후 $B_{\rm rel}>10^{-3}$ 경계 근방 꼬리 정밀화. $N_{\rm null}=1000$ 재평가 후 네 사례 모두 엄격한 국소 최댓값 기준 $PTE_{\rm global,max}<0.01$을 통과하지 못한다.}
\label{tab:tail_refined_cases}
\begin{tabularx}{\textwidth}{>{\raggedright\arraybackslash}p{3.2cm}
>{\centering\arraybackslash}p{2.4cm}
>{\centering\arraybackslash}p{1.9cm}
>{\centering\arraybackslash}p{1.9cm}
>{\centering\arraybackslash}p{1.8cm}
>{\centering\arraybackslash}p{1.6cm}}
\toprule
셀/조건 & 형태학 & $PTE_{\rm max}^{100}$ & $PTE_{\rm max}^{1000}$ & $PTE_{\rm absint}^{1000}$ & $B_{\rm rel}$ \\
\midrule
$0\_197$ 빈 36 전체 빔/다중 & 매끄러운 광대역 1000 Jy & 0.0198 & 0.0220 & 0.913 & 0.0132 \\
$0\_197$ 빈 36 전체 빔/다중 & kHz-comb 1000 Jy & 0.0297 & 0.0250 & 0.903 & 0.0117 \\
$0\_197$ 빈 0 고정 빔/단일 & 매끄러운 광대역 1000 Jy & 0.0297 & 0.0400 & 0.0889 & 0.0096 \\
$23\_80$ 빈 48 전체 빔/단일 & 협대역 선형 1000 Jy & 0.0198 & 0.0340 & 0.969 & 0.0018 \\
\bottomrule
\end{tabularx}
\end{table*}

편향 기준값 민감도도 후보 해석에 중요하다. $B_{\rm floor}=10^{-3}$에서는
\BiasFloorOneEThreeCandidates{}건, $10^{-2}$에서는 \BiasFloorOneETwoCandidates{}건, $10^{-1}$에서는
\BiasFloorOneEOneCandidates{}건이 꼬리 정밀화 후 남는다. 따라서 국소 최댓값 PTE 후보 수는 절대 편향 기준값의
정의에 의존하며, PTE 단독으로 물리적 오염을 주장할 수 없다. 본 연구는 $B_{\rm floor}=10^{-2}$를
보편 권고값이 아니라 본 시험 영역의 보고 수준 진폭 게이트로만 사용한다. 생산 수준 분석에서는
열잡음 공분산, 보정 잔차 규모, 전경 누출 예산 등으로 사전에 정한 물리적 기준값이
이 역할을 대체해야 한다.

\begin{figure*}[!t]
\centering
\includegraphics[width=0.82\textwidth]{R3_bias_floor_sensitivity.png}
\caption{절대 편향 기준값 민감도. 엄격한 국소 최댓값 후보 수는 $B_{\rm floor}$ 설정에 따라 달라질 수 있다. $B_{\rm floor}=10^{-2}$는 본 시험 영역의 보고 수준 진폭 게이트이며, 보편적인 물리 상수가 아니다. 실제 생산 수준 QA에서는 해당 관측의 잡음 및 계통오차 예산에 근거하여 $B_{\rm floor}$를 사전에 정해야 한다.}
\label{fig:bias-floor}
\end{figure*}

\subsection{윈도우 적분 PTE: 빔-강건 광역 밴드파워 초과는 없다}
\label{sec:results:absint}

꼬리 정밀화된 경계 근방 사례 \NTailRefinedCases{}건은 모두 $PTE_{\rm global,max}<0.01$을 만족하지 못한다.
또한 이 사례들의 윈도우 적분 통계량 $PTE_{\rm global,absint}$는 \PTEabsintStrictRange{} 범위에 머문다.

한편 $PTE_{\rm global,absint}$ 특화 바닥값 재검증에서는 다른 결과가 관측되었다. 거친 그리드에서
$PTE_{\rm global,absint}$ 바닥값($\PTEMinHundred$)에 도달한 $B_{\rm rel}>10^{-3}$인 \NAbsintFloorAboveGate{}건은
$N_{\rm null}=1000$ 재평가에서도 $PTE_{\rm global,absint}<0.01$이 유지되었다(고정 PolyBeam/매끄러운/1000 Jy:
$PTE_{\rm absint}=0.006$, $n_{\rm exceed}=5/1000$; 고정 PolyBeam/kHz-comb/1000 Jy: $PTE_{\rm absint}=0.008$,
$n_{\rm exceed}=7/1000$). 이 초과는 $PTE_{\rm global,max}\approx0.74$에서도 나타나지 않는 분산된 적분
초과로, 두 위성의 코히런트 합산이 위성별 전역 위상 무작위화 대조군보다 더 많은 윈도우 적분 파워를
남기는 패턴에 기인한다.

그러나 이 두 사례는 모두 고정 PolyBeam에서만 관찰되며, 짝지은 전체 PolyBeam 시행의
$PTE_{\rm global,absint}$는 각각 0.030과 0.020으로, $N=100$ 단계에서도 이미 엄격 기준을 통과하지
못한다. 따라서 빔-강건 윈도우 적분 후보는 \NAbsintBeamRobust{}건이다. 수정된 커버리지 검정에서는
국소 최댓값과 빔-강건 윈도우 적분 양쪽 모두에서 밴드파워 오염 후보가 확인되지 않는다.

\subsection{전체 PolyBeam 강건성: 빔-강건 후보는 없다}
\label{sec:results:beam_robustness}

수정된 꼬리 정밀화에서는 고정 PolyBeam 엄격 후보도, 전체 PolyBeam 엄격 후보도 남지 않는다. 따라서
빔-강건 엄격 후보는 \NStrictBeamRobust{}건이며, 전체 PolyBeam 엄격 후보도 \NFullPolyStrict{}건이다.
초기 빈 중심 프록시 결과에서 보였던 외관상의 단일 셀 반응은 빈 전체 위성 메타데이터 수정 후
선택된 셀 자체가 바뀌면서 유지되지 않았다. 이는 LST 빈 층화가 위성 통과 기하를 충분히
반영해야 하며, 빈 중심 노출 프록시만으로 ``정온'' 또는 ``전형'' 분류를 물리적으로 해석하면 안 된다는
방법론적 경고를 제공한다.

\begin{table*}[!t]
\centering
\small
\caption{수정된 꼬리 정밀화 요약. 빈 전체 위성 노출 메타데이터를 사용한 재실행 후 엄격한 국소 최댓값 후보는 남지 않는다. 하단 블록은 $PTE_{\rm global,absint}$ 특화 바닥값($N=100$) 사례의 별도 $N=1000$ 재검증 결과이며, 빔-강건 기준에서는 후보가 없다.}
\label{tab:tail_summary}
\begin{tabularx}{\textwidth}{X >{\centering\arraybackslash}p{2.4cm}}
\toprule
지표 & 값 \\
\midrule
\multicolumn{2}{l}{\textit{$PTE_{\rm global,max}$ 꼬리 정밀화}} \\
정밀화된 경계 근방 사례 & \NTailRefinedCases \\
$PTE_{\rm global,max}<0.01$ & \NStrictLocalMax \\
$PTE_{\rm global,max}<0.01$ 및 $B_{\rm rel}>10^{-3}$ & \BiasFloorOneEThreeCandidates \\
$PTE_{\rm global,max}<0.01$ 및 $B_{\rm rel}>10^{-2}$ & \BiasFloorOneETwoCandidates \\
$PTE_{\rm global,max}<0.01$ 및 $B_{\rm rel}>10^{-1}$ & \BiasFloorOneEOneCandidates \\
전체 PolyBeam 엄격 후보 & \NFullPolyStrict \\
빔-강건 엄격 후보 & \NStrictBeamRobust \\
고정 빔 전용 엄격 후보 & \NFrozenOnlyStrict \\
\midrule
\multicolumn{2}{l}{\textit{$PTE_{\rm global,absint}$ 바닥값 재검증 ($N=1000$)}} \\
거친 그리드 absint 바닥값 사례 & \NAbsintFloorHitsCoarse \\
$B_{\rm rel}>10^{-3}$ 재검증 사례 & \NAbsintFloorAboveGate \\
$N=1000$ 후 $PTE_{\rm absint}<0.01$ 유지 (고정 빔 전용) & \NAbsintRetainedN1000 \\
빔-강건 윈도우 적분 후보 & \NAbsintBeamRobust \\
\bottomrule
\end{tabularx}
\end{table*}

\FloatBarrier

\section{논의}
\label{sec:discussion}

본 연구의 결과는 다섯 가지 점을 강조한다.

첫째, 국소 강건 z 또는 국소 최댓값 PTE는 오염 판정 기준이 아니라 QA 진단량이다.
$Z_{\rm PS,max}$는 윈도우 내 최댓값 통계량이며, 분모가 대응-영가설 MAD로 정규화되므로 영가설 분산이
작은 조건에서 절대 편향과 분리된 국소 플래그를 만들 수 있다. 따라서 $Z_{\rm PS,max}>3$ 또는
$PTE_{\rm global,max}<0.01$을 곧바로 EoR 윈도우 오염 검출로 해석하면 안 된다.

둘째, 국소 최댓값 초과와 윈도우 적분 밴드파워 오염은 분리되어야 하며, 두 종류의 PTE는 서로 독립적으로 점검해야
한다. $PTE_{\rm global,max}$ 기준 꼬리 정밀화 대상 \NTailRefinedCases{}건은 $N_{\rm null}=1000$ 재평가 후
$PTE_{\rm global,max}<0.01$을 만족하지 못했고, 이 사례들의 $PTE_{\rm global,absint}$도 0.0889 이상이었다.
한편 $PTE_{\rm global,absint}$ 특화 바닥값 사례 중 $B_{\rm rel}>10^{-3}$인 \NAbsintFloorAboveGate{}건은
$N_{\rm null}=1000$에서도 $PTE_{\rm global,absint}<0.01$이 유지되었으나, 짝지은 전체 PolyBeam 시행에서
기준을 통과하지 못해 빔-강건 후보에는 해당하지 않는다. 이는 $PTE_{\rm global,absint}$ 바닥값 도달 사례도
$PTE_{\rm global,max}$ 기준 꼬리 정밀화와 별도로 점검되어야 함을 보여주며, 두 통계량을 단일 선별 기준으로
묶어서는 안 된다는 방법론적 경고를 제공한다.

셋째, 빔 연산자 강건성은 필수적인 판정 단계이다. 수정된 꼬리 정밀화에서는 고정 PolyBeam 엄격
후보도 전체 PolyBeam 엄격 후보도 남지 않았다. 따라서 본 연구의 최종 결과는 빔-강건 오염 검출이
아니라, 수정된 층화와 다단계 판정 규칙을 통과하는 후보가 없다는 비검출이다.

넷째, LST 빈 층화 메타데이터는 빈 중심 프록시가 아니라 native 10분 청크 전체의 위성 기하를
반영해야 한다. 초기 점검에서 $(0\_1,\mathrm{bin}\ 36)$은 청크 중심 기준 위성 노출이 0이었지만,
청크 전체 시간 표본 순위에서는 고고도 위성 통과를 포함하였다. 이 불일치를 9개 기선 전체의
메타데이터 산출에 적용해 수정하자 선택된 셀이 27개 중 15개 항목에서 바뀌었고, 기존 외관상의
단일 셀 엄격 패턴은 유지되지 않았다. 따라서 ``전형'' 분류는 사용된 메타데이터 연산자에 조건부인 분류이며,
위성 스트레스가 없는 일반 상태라는 물리적 해석으로 곧장 사용해서는 안 된다.

다섯째, 본 결과는 Starlink 유사 UEMR의 일반적 안전성을 주장하지 않는다. 본 연구는 시험된 주입 영역과
단순화된 HERA 유사 QA 파이프라인에서 빔-강건 윈도우 적분 밴드파워 오염을 검출하지 못했다는
제한된 음성 결과와, 빈 전체 메타데이터 수정이 결과 해석을 바꿀 수 있다는 방법론적 경고를 함께 제시한다.
실제 생산 수준 HERA 검증은 전체 $k_\perp$--$k_\parallel$ 전파,
열잡음 공분산, 중복도 가중치, LST 빈 묶음, 신호 손실 보정, 보정/계통오차
모델링을 포함해야 한다.

마지막으로, 본 연구의 새로움은 Starlink UEMR 자체의 최초 보고가 아니라 QA 판정 계층의 명시적 분리이다.
Di Vruno 등과 Bassa 등 계열 연구는 저주파 배열에서 Starlink UEMR의 존재와 스펙트럼 구조를 관측적으로 제시한다
\citep{divruno2023,bassa2024}. 반면 본 연구는 그러한 문헌 기반 UEMR 구조가 HERA 유사 지연 스펙트럼 QA에서
어떤 통계량으로 남는지를 묻는다. HERA 상한 및 잭나이프/영가설 검정 계열 연구는 밴드파워, PTE, 종단간 모의실험,
계통오차 모델링을 결합하여 최종 우주론 산출물을 검증한다 \citep{abdurashidova2022,wilensky2019}. 본 연구의
다단계 규칙은 그 생산 수준 검증을 대체하지 않고, 그 이전 단계에서 국소 최댓값 초과를 윈도우 적분 밴드파워 및
빔-강건 오염 주장으로 과대해석하지 않도록 분류하는 생산 전(pre-production) QA 계층으로 위치한다.

\subsection{제한점}
\label{sec:discussion:limitations}

본 연구의 가장 중요한 제한점은 배경 실현 분산을 충분히 표본화하지 못했다는 점이다. 커버리지 검정은
2차원 층화추출로 고른 9개 native HERA H6C EW-like 기선을 사용하지만, 이들은 같은 관측/처리 계열과
동일한 Jan 2026 TLE 재매핑 가정 아래 평가된다. 따라서 초록과 결론의 음성 결과는 반드시 이 배경 실현,
단순화된 지연 스펙트럼 QA 연산자, 선택된 LST 커버리지, 그리고 사용한 빔 모델에 조건부이다. 특히 ``다른 HERA 관측 시즌, 다른 LST 범위, 다른 플래깅
상태에서도 동일한 후보 수가 나오는가''라는 질문에는 현재 실험만으로 답할 수 없다. 본 논문의 음성 결과는 Starlink 유사 UEMR의
일반적 안전성 주장이 아니라, 주어진 조건에서 빔-강건 윈도우 적분 오염이 검출되지 않았다는 제한된 QA 결과로 해석해야 한다.

두 번째 제한점은 시행들이 독립적인 우주론 실현이 아니라는 점이다. 기선, 형태학, 플럭스, 빔 연산자, LST 빈은
동일한 배경과 처리 가정을 공유한다. 따라서 $648\times0.01$ 기대 가짜 양성 계산은 엄밀한 가족 단위
오류율(family-wise error rate) 보정이 아니다. 이 계산은 전체 후보 수의 규모를 확인하는 보조 참고값으로만 사용한다.
향후 분석에서는 전 그리드 최댓값 통계량, 순열 기반 보정, 또는 계층적 배경 모델을 사용해 그리드 전체의
오경보 확률을 직접 보정할 필요가 있다.

세 번째 제한점은 대응-영가설의 구현 방식과 실험별 해석 차이이다. 통제 그리드에서는 적분 시간당 하나의 독립 위상을
적용하는 per-time 위상 스크램블을 사용하므로 스펙트럼 결맞음(지연 스파이크)이 보존된다. 커버리지 검정에서는
위성별 하나의 전역 위상 스칼라를 적용하는 per-satellite 위상 무작위화를 사용하므로, 각 위성의 지연--시간
코히런트 구조 전체가 영가설에서도 보존된다. 두 경우 모두, 영가설은 ``UEMR 신호가 없다''는 영가설이 아니라
``결맞은 조직화가 위상 무작위화 대조군보다 얼마나 드문가''를 측정하는 조건부 스트레스 테스트이다. 특히 커버리지
단일 위성 시행의 영가설에서는 배경--위성 교차항 위상만 달라지며, $PTE_{\rm global}$는 결맞음 교차항 기여가
위상 무작위화 대조군보다 통계적으로 극단적인지를 묻는 지표이다.
이 구분은 본 연구의 비검출 서술과 판정 계층 해석에 필수적이다.

네 번째 제한점은 절대 편향 게이트의 물리적 캘리브레이션이다. 본 연구에서 보고된 $B_{\rm floor}$ 값들은 이 단순화된 지연 스펙트럼 QA 연산자 내에서 상대적인 EoR 윈도우 편향에 대한 무차원 임계값이다. 즉, $\mathrm{mK}^2$ 단위의 직접적인 HERA 우주론적 상한 임계값이 아니다. $B_{\rm floor}=10^{-2}$와 같은 값을 실제 파워 스펙트럼 허용 오차에 매핑하려면, 대상 분석에 대한 전체 밴드파워 정규화, 열잡음 공분산, 신호 손실 보정, 중복도 가중치, 그리고 보정 및 계통오차 예산이 필요하다. 운영 측면에서 이 임계값은 부분적인 윈도우 오염에 대한 보고 수준의 게이트로 해석되어야 하며, 실제 HERA 파이프라인에서는 유도된 편향이 각 $k_\perp$--$k_\parallel$ 빈에서 사전 지정된 밴드파워 오차 예산보다 하위(subdominant) 수준이 되도록 설정되어야 한다.

다섯 번째 제한점은 kHz-comb 형태학에서 Doppler 편이 효과의 미반영이다. 본 주입 모델은 위성의 기하학적
지연 $\tau_{\rm sat}(t)$를 결정론적으로 계산하지만, 시선 방향 거리 변화율(range rate $\dot{r}$)에서
비롯되는 Doppler 주파수 편이는 스펙트럼에 반영하지 않는다. LEO Starlink 위성의 전형적 시선
속도($\sim$7 km/s)는 150\,MHz 대역에서 $\dot{r}/c\,\nu \approx 3.5$\,kHz 수준의 편이를 만든다. 10분
통과 동안 $\dot{r}$은 수 km/s 변화하므로 kHz-comb 라인은 관측 구간 내에서 수 kHz 드리프트하며,
이는 kHz-comb 간격과 동급의 주파수 이동이다. 따라서 현재 주입은 kHz-comb 형태학의 시간적 위상
결맞음을 과대평가할 수 있으며, 실제 Doppler drift가 적용되면 comb 라인의 시간 방향 탈코히런스로
인해 윈도우 기여가 감소할 가능성이 있다. 이 효과는 향후 주입 모델에 Doppler 추적을 추가하여
정량화해야 한다.

\section{결론}
\label{sec:conclusion}

본 원고는 다음의 결론을 지지한다.
\begin{enumerate}
  \item Starlink 유사 UEMR의 존재와 EoR 윈도우 밴드파워 오염은 동일하지 않으며, 처리 후 잔차는 국소, 윈도우 적분, 빔-강건 단계로 분리되어 평가되어야 한다.
  \item 통제 그리드에서 국소 강건 z 플래그는 절대 편향 및 전역 PTE와 일대일로 대응하지 않았으며, $Z_{\rm PS,max}$ 단독은 오염 판정 기준으로 부적절하다.
  \item LST 빈 위성 노출 메타데이터를 빈 전체 시간 표본 기준으로 수정하자 27개 선택된 층 항목 중 15개가 바뀌었고, 기존 외관상의 단일 셀 엄격 패턴은 유지되지 않았다.
  \item 11시간 기선--LST 커버리지 검정에서 $B_{\rm rel}>10^{-3}$ 경계 근방 \NTailRefinedCases{}건을 $N_{\rm null}=1000$으로 재평가한 결과, 꼬리 정밀화된 엄격한 국소 최댓값 후보는 \NStrictLocalMax{}건이었다.
  \item $PTE_{\rm global,absint}$ 특화 바닥값 재검증에서 $B_{\rm rel}>10^{-3}$인 \NAbsintRetainedN1000{}건이 $N_{\rm null}=1000$에서도 $PTE_{\rm global,absint}<0.01$로 확인되었으나, 짝지은 전체 PolyBeam 시행의 $PTE_{\rm global,absint}\ge0.02$로 빔-강건 기준을 통과하지 못하였다. 따라서 빔-강건 윈도우 적분 후보는 \NAbsintBeamRobust{}건이다.
  \item 짝지은 전체 PolyBeam 엄격 후보는 \NFullPolyStrict{}건이고, 빔-강건 후보도 \NStrictBeamRobust{}건이었다.
  \item 따라서 본 연구는 시험된 영역에서 Starlink 유사 UEMR의 빔-강건 윈도우 적분 EoR 윈도우 오염을 검출했다고 주장하지 않는다.
  \item 본 연구의 방법론적 기여는 국소 최댓값 PTE, 절대 편향 기준값, 윈도우 적분 PTE, 전체 빔 강건성을 분리한 다단계 UEMR QA 판정 규칙을 제안하는 데 있다.
\end{enumerate}

요약하면, 본 연구는 Starlink 유사 UEMR이 곧바로 21-cm 파워 스펙트럼을 오염시킨다고 주장하지 않는다.
시험된 HERA H6C native 9-baseline 조건과 주입 영역 안에서 빔-강건 윈도우 적분 오염은 검출되지 않았고,
수정된 꼬리 정밀화 후 엄격한 국소 최댓값 후보도 남지 않았다.
따라서 본 연구의 기여는 빈 전체 위성 메타데이터, 꼬리 정밀화된 PTE, 절대 편향 기준값, 윈도우 적분 통계량,
전체 빔 강건성을 함께 요구하는 다단계 QA 구조를 제시하는 데 있다.

\section*{데이터 및 코드 가용성}
재현성 코드, 설정 파일, 주입 실험 산출물 및 본 원고에서 인용한 결과 표는 GitHub 저장소
\url{https://github.com/kimeujin03-droid/starlink_uemr}에서 제공한다. 제출 또는 출판 시점에는 해당 공개판을 Zenodo에
아카이브하여 DOI를 부여하고, 원고에는 GitHub URL과 버전이 지정된 DOI를 함께 인용한다. 분석에 사용한 Python, NumPy, SciPy,
pandas, Astropy, pyuvdata, healpy 및 hera\_sim의 정확한 버전, 운영체제, 난수 시드 정책, 그리고 $N_{\rm null}=100/1000$
설정은 저장소의 \texttt{environment.yml}, \texttt{requirements.txt}, \texttt{REPRODUCIBILITY.md} 및 각 \texttt{config\_used.yaml}에 기록한다.
원본 HERA 유사 배경 가시성이 재배포 불가능한 경우, 본문 그림과 표를 재생성할 수 있는 요약 CSV와 그림 생성 스크립트를 함께 제공한다.

\section*{감사의 글}
본 연구는 공개 HERA 파생 가시성 산출물과 순방향 주입 워크플로우를 위해 개발된 로컬 재현성 스크립트를 사용하였다.

\bibliographystyle{elsarticle-harv}
\bibliography{starlink_uemr_refs,starlink_uemr_stat_refs_additions}

\end{document}
